\chapter{Graph Theory --- Basic}
%%% generated by the blueprint pipeline from TCSlib.GraphTheory.Kruskal.Basic
\section{Overview}

This module formalises the core data structures and the greedy selection loop
underlying Kruskal's minimum-spanning-tree algorithm.  It defines weighted edges
over $n$ vertices, a union-find structure represented as a flat array of
representatives, the edge-processing sweep that builds the spanning tree, and the
top-level \texttt{kruskal} function.  The two main correctness results show that
every edge returned by the algorithm was present in the original input list.

%%%%%%%%%%%%%%%%%%%%%%%%%%%%%%%%%%%%%%%%%%%%%%%%%%%%%%%%%%%%%%%%%%%%%%%%%%%%%%%
\section{Declarations}

\begin{definition}[Weighted edge]
\lean{Kruskal.WEdge}
\label{Kruskal.WEdge}
\leanok
A \emph{weighted edge} over $n$ vertices is a structure with two endpoints
$u, v : \mathrm{Fin}\,n$ and a natural-number weight $w : \mathbb{N}$.
\end{definition}

\begin{definition}[Union-find type]
\lean{Kruskal.UF}
\label{Kruskal.UF}
\leanok
The union-find state for $n$ nodes is defined as a function
$\mathrm{Fin}\,n \to \mathbb{N}$ mapping each node to its current
representative (component label).
\end{definition}

\begin{definition}[Initial union-find]
\lean{Kruskal.UF.init}
\label{Kruskal.UF.init}
\leanok
\uses{Kruskal.UF}
$\mathtt{init}\,n$ constructs the identity union-find on $n$ nodes: every node
$i$ is its own representative, i.e.\ $\mathtt{init}\,n\,(i) = i$.
\end{definition}

\begin{definition}[Find representative]
\lean{Kruskal.UF.find}
\label{Kruskal.UF.find}
\leanok
\uses{Kruskal.UF}
$\mathtt{find}\,\mathit{uf}\,i$ returns the representative of node $i$ in
the union-find state $\mathit{uf}$, defined simply as $\mathit{uf}(i)$.
\end{definition}

\begin{definition}[Merge two components]
\lean{Kruskal.UF.merge}
\label{Kruskal.UF.merge}
\leanok
\uses{Kruskal.UF}
$\mathtt{merge}\,\mathit{uf}\,i\,j$ unifies the components of nodes $i$ and
$j$ by reassigning every node whose representative equals $\mathit{uf}(j)$ to
instead carry the representative $\mathit{uf}(i)$.
\end{definition}

\begin{lemma}[Find on the initial union-find]
\lean{Kruskal.UF.init_find}
\label{Kruskal.UF.init_find}
\leanok
\uses{Kruskal.UF.find, Kruskal.UF.init}
\difficulty{1}
Consider the identity union-find on $n$ nodes, in which every node is its own
representative. For every node $i \in \mathrm{Fin}\,n$, applying the find operation to
this initial state returns the underlying numeric value of $i$; that is, the
representative of $i$ is $i$ itself.
\end{lemma}

\begin{lemma}[Find equals direct lookup]
\lean{Kruskal.UF.find_def}
\label{Kruskal.UF.find_def}
\leanok
\uses{Kruskal.UF, Kruskal.UF.find}
\difficulty{1}
Let $\mathit{uf} \colon \mathrm{Fin}\,n \to \bbn$ be a union-find state on $n$ nodes,
and let $i$ be one of its nodes. Then the representative of $i$ returned by
$\mathrm{find}$ coincides with the direct evaluation of the state at $i$; that is,
$\mathrm{find}(\mathit{uf}, i) = \mathit{uf}(i)$.
\end{lemma}

\begin{definition}[Merge all edges into union-find]
\lean{Kruskal.UF.mergeAll}
\label{Kruskal.UF.mergeAll}
\leanok
\uses{Kruskal.UF, Kruskal.UF.merge, Kruskal.WEdge}
$\mathtt{mergeAll}\,\mathit{uf}\,\mathit{es}$ folds a list of weighted edges
$\mathit{es}$ into the union-find $\mathit{uf}$ by calling \texttt{merge} on
each edge's endpoints in sequence.
\end{definition}

\begin{definition}[Same partition predicate]
\lean{Kruskal.UF.SamePartition}
\label{Kruskal.UF.SamePartition}
\leanok
\uses{Kruskal.UF}
$\mathtt{SamePartition}\,\mathit{uf}_1\,\mathit{uf}_2$ holds when two
union-find states induce the same equivalence relation on nodes: for all
$i, j : \mathrm{Fin}\,n$, $\mathit{uf}_1(i) = \mathit{uf}_1(j)$ if and
only if $\mathit{uf}_2(i) = \mathit{uf}_2(j)$.
\end{definition}

\begin{definition}[Edge-selection sweep]
\lean{Kruskal.processEdges}
\label{Kruskal.processEdges}
\leanok
\uses{Kruskal.UF, Kruskal.UF.find, Kruskal.UF.merge, Kruskal.WEdge}
$\mathtt{processEdges}\,\mathit{es}\,\mathit{uf}\,\mathit{acc}$ scans the
edge list $\mathit{es}$ left-to-right: an edge $e$ is added to the
accumulator and its endpoints are merged whenever
$\mathit{uf}.\mathtt{find}\,e.u \ne \mathit{uf}.\mathtt{find}\,e.v$;
otherwise $e$ is skipped.  The reversed accumulator is returned when the
list is exhausted.
\end{definition}

\begin{definition}[Kruskal's algorithm]
\lean{Kruskal.kruskal}
\label{Kruskal.kruskal}
\leanok
\uses{Kruskal.UF.init, Kruskal.WEdge, Kruskal.processEdges}
$\mathtt{kruskal}\,n\,\mathit{edges}$ sorts the edge list by weight and then
runs \texttt{processEdges} on the sorted list starting from the identity
union-find, returning the selected spanning-tree edges.
\end{definition}

\begin{lemma}[Membership in the edge-selection output]
\lean{Kruskal.processEdges_mem}
\label{Kruskal.processEdges_mem}
\leanok
\uses{Kruskal.UF, Kruskal.WEdge, Kruskal.processEdges}
\difficulty{4}
Fix $n$, and let $\mathit{es}$ be a list of weighted edges over $n$ vertices, let
$\mathit{uf}$ be a union-find state on $n$ nodes, and let $\mathit{acc}$ be an
accumulator list of weighted edges. If a weighted edge $e$ occurs in the output produced
by running the edge-selection sweep on the edge list $\mathit{es}$, the union-find state
$\mathit{uf}$, and the accumulator $\mathit{acc}$, then $e$ already occurs in the
accumulator $\mathit{acc}$ or $e$ occurs in the edge list $\mathit{es}$.
\end{lemma}

\begin{lemma}[Kruskal's algorithm returns a subset of its input edges]
\lean{Kruskal.kruskal_subset}
\label{Kruskal.kruskal_subset}
\leanok
\uses{Kruskal.WEdge, Kruskal.kruskal, Kruskal.processEdges_mem}
\difficulty{3}
Fix $n \in \bbn$ and a list of weighted edges over $n$ vertices. Every edge returned by
running Kruskal's algorithm on this list occurs in the original list: if $e$ belongs to
the output of $\mathtt{kruskal}\,n\,\mathit{edges}$, then $e$ belongs to
$\mathit{edges}$.
\end{lemma}
